\documentclass[11pt]{article}
\usepackage[margin=1in]{geometry}
\usepackage{amsmath,amssymb,amsthm,booktabs,hyperref}
\newtheorem{theorem}{Theorem}\newtheorem{lemma}[theorem]{Lemma}\newtheorem{conjecture}[theorem]{Conjecture}
\theoremstyle{remark}\newtheorem{remark}[theorem]{Remark}
\title{Rusty numbers are linear plus automatic\\ for \(d\le 8\) and \(d=11\)}
\author{Aditya Balaji}
\date{24 November 2026}
\begin{document}\maketitle
\begin{abstract}
Fokkink and Joshi asked whether the \emph{rusty numbers}, the anti-recurrence sequences for the linear form \((1,d)\), satisfy the Clergyman's conjecture: is \(A_n-(2d+3)n\) a \((d+1)\)-automatic sequence? The cases \(d=1,2\) were known. We prove the conjecture for \(d=3,4,5,6,7,8\) and \(d=11\) by exhibiting explicit automata and checking the defining recursion with the theorem prover Walnut. The automaton for \(d=4\) explains the break in the progression \(A_{5n+1}=9+55n\) at \(n=348\).
\end{abstract}

\section{The problem}
Fix \(d\ge1\). The rusty sequence \(A=(A_n)_{n\ge1}\) and its companion \(B=(B_m)_{m\ge1}\) are built greedily \cite{fj}. For \(n=1,2,\dots\):
\begin{itemize}
\item \(B_{2n-1}\) is the least positive integer not among \(A_1,\dots,A_{n-1},B_1,\dots,B_{2n-2}\);
\item \(B_{2n}\) is the least integer greater than \(B_{2n-1}\) not among \(A_1,\dots,A_{n-1}\);
\item \(A_n=B_{2n-1}+d\,B_{2n}\).
\end{itemize}
Thus \(A\) and \(B\) partition the positive integers, and \(A_1=1+2d\). The trace is \(\tau=1+d\) and \(\kappa=2\tau+1=2d+3\). The Clergyman's conjecture of Kimberling and Moses, in the form of \cite{fj}, predicts that \(A_n-\kappa n\) is \(\tau\)-automatic. Fokkink and Joshi prove it for every \(A_1\)-bounded form, which covers \(d\le2\). They note that \((1,d)\) is not \(A_1\)-bounded for \(d\ge3\), call the rusty numbers ``a challenge to prove, or disprove'', and observe that for \(d=4\) the progression \(A_{5n+1}=9+55n\) fails first at \(n=348\).

\begin{theorem}\label{main}
For \(d\in\{1,2,3,4,5,6,7,8,11\}\), the sequence \(A_n-(2d+3)n\) is \((d+1)\)-automatic. It is generated by the msd-first deterministic finite automaton with output listed in Table~\ref{tab}.
\end{theorem}

\begin{table}[h]\centering
\begin{tabular}{rrrrl}\toprule
\(d\) & \(\kappa\) & base & states & status\\\midrule
1&5&2&2&known (anti-Fibonacci)\\
2&7&3&3&known (anti-Pell)\\
3&9&4&6&new\\
4&11&5&9&new\\
5&13&6&11&new\\
6&15&7&16&new\\
7&17&8&19&new\\
8&19&9&21&new\\
11&25&12&31&new\\\bottomrule
\end{tabular}
\caption{Automata for \(A_n-\kappa n\) (\texttt{RUSTY*.txt} in the repository).}\label{tab}
\end{table}

\section{Proof}
For each \(d\), let \(F\) be the automaton of Table~\ref{tab} and define candidate sequences
\[
a_n=\kappa n+F(n-1)-c_d,\qquad b^{(2)}_n=\Bigl\lfloor\frac{a_n+2}{d+1}\Bigr\rfloor,\qquad b^{(1)}_n=a_n-d\,b^{(2)}_n,
\]
with a constant offset \(c_d\) fixed by the automaton encoding. Each of these is synchronized in base \(d+1\), so first-order statements about them are decidable \cite{shallit}. Walnut \cite{walnut} verified the following eight statements for every \(d\) in the table, printing \texttt{TRUE} each time.
\begin{enumerate}
\item[(T)] \(a_n,b^{(1)}_n,b^{(2)}_n\) are defined for all \(n\ge1\).
\item[(I)] \(b^{(1)}_1=1\), \(b^{(2)}_1=2\), \(a_1=1+2d\).
\item[(O)] \(b^{(1)}_n<b^{(2)}_n<b^{(1)}_{n+1}\) for all \(n\).
\item[(N1)] \(b^{(1)}_n\ne a_m\) for \(0<m<n\).
\item[(N2)] \(b^{(2)}_n\ne a_m\) for \(0<m<n\).
\item[(M1)] every \(t\) with \(0<t<b^{(1)}_n\) equals \(a_m\), \(b^{(1)}_m\) or \(b^{(2)}_m\) for some \(0<m<n\).
\item[(M2)] every \(t\) with \(b^{(1)}_n<t<b^{(2)}_n\) equals \(a_m\) for some \(0<m<n\).
\item[(S)] \(a_n=b^{(1)}_n+d\,b^{(2)}_n\).
\end{enumerate}

\begin{lemma}
If candidate sequences satisfy (T)--(S), then \(a_n=A_n\), \(b^{(1)}_n=B_{2n-1}\) and \(b^{(2)}_n=B_{2n}\) for all \(n\).
\end{lemma}
\begin{proof}
Induction on \(n\); assume agreement for all \(m<n\). By (O), \(b^{(1)}_n\) exceeds every earlier \(b\), and by (N1) it is no earlier \(a\), so it is unused; by (M1) every smaller positive integer is used. Hence \(b^{(1)}_n=B_{2n-1}\). Likewise \(b^{(2)}_n>b^{(1)}_n\) is no earlier \(a\) by (N2), and by (M2) every integer strictly between is an earlier \(a\). Hence \(b^{(2)}_n=B_{2n}\), and (S) gives \(a_n=A_n\).
\end{proof}
This proves Theorem~\ref{main}; (I) is implied but serves as a sanity check. The same eight-statement method is used in \cite[Theorem~5]{fj}.

\begin{remark}
For \(d=4\) one state of the automaton is first reached at prefix \(348\). Before that point the automaton's output agrees with the arithmetic progression; this is the break observed in \cite{fj}.
\end{remark}

\section{How the automata were found}
A C program generated \(10^8\)--\(4\times10^8\) terms of \(A_n-\kappa n\), cross-checked against an independent Python implementation on \(2\times10^6\) terms. The msd automaton was built by Myhill--Nerode merging of prefixes and validated against all generated terms. This search is heuristic, but the proof above does not depend on it.

For \(d=1,\dots,30\) the data suggest that \(A_n-\kappa n\) is bounded, with values in roughly \([-1.3d,\,0.57d]\), and is the fixed point of a \(k\)-block \((d+1)\)-uniform substitution, with \(k=1\) for \(d\le9\), \(k=2\) for \(10\le d\le13\), \(k=3\) for \(d=14,15\), and \(k=5,6,6\) for \(d=20,25,30\). For \(d=9,10\) and \(d\ge12\), \(10^8\) terms did not suffice to extract an automaton because some states are first reached very late.

\begin{conjecture}
For every \(d\ge1\), \(A_n-(2d+3)n\) is bounded and \((d+1)\)-automatic.
\end{conjecture}

\section{Reproducibility and disclosure}
Walnut 6 was compiled from source. All runs used at most \(0.71\)\,GB of memory, except \(d=11\), which peaked at \(2.7\)\,GB (Java heap limit \(4\)\,GB). The automata, Walnut command files, logs and generators are included with this note. The computations, automaton search and drafting were done with AI assistance (Anthropic's Claude). The author checked the results and is responsible for the content.

\begin{thebibliography}{9}
\bibitem{fj} R. Fokkink and G. Joshi, \emph{Anti-recurrence sequences}, arXiv:2506.13337, 2025.
\bibitem{walnut} H. Mousavi, \emph{Automatic theorem proving in Walnut}, arXiv:1603.06017, 2016.
\bibitem{shallit} J. Shallit, \emph{The Logical Approach to Automatic Sequences: Exploring Combinatorics on Words with Walnut}, London Math. Soc. Lecture Note Ser. 482, Cambridge Univ. Press, 2022.
\end{thebibliography}
\end{document}
